\usepackage{fontspec,xunicode,xltxtra}
% \setmainfont[Mapping=tex-text]{Georgia}
\setmainfont[Mapping=tex-text]{Arial}
\setsansfont[Mapping=tex-text]{Arial}
\setmonofont[Mapping=tex-text]{Courier New}
% \setmonofont[Mapping=tex-text]{Verdana}
% \setmonofont[Mapping=tex-text]{Consolas}

% \usepackage{xeCJK}
% \setCJKmainfont[ItalicFont={楷体_GB2312}]{宋体}
% \setCJKsansfont{黑体}
% \setCJKmonofont{黑体}

\usepackage{xeCJK}
% \setCJKmainfont[ItalicFont={Adobe Kaiti Std}]{Adobe Song Std}
\setCJKmainfont[ItalicFont={Adobe Kaiti Std}]{Adobe Heiti Std}
\setCJKsansfont{Adobe Heiti Std}
% \setCJKsansfont{Microsoft YaHei}
\setCJKmonofont{Adobe Kaiti Std}
\punctstyle{banjiao}

\graphicspath{{../fig/}}
% \usepackage{indentfirst}

\usepackage[dvipdfmx]{hyperref}  % options: dvipdfmx, pdftex 
\hypersetup{
    % bookmarks=true,         % show bookmarks bar?
    bookmarksopen=false,
    pdfpagemode=UseNone,    % options: UseNone, UseThumbs, UseOutlines, FullScreen
    pdfstartview=FitH,      % options: FitH, FitV
    pdfborder=1,
    pdfhighlight=/P,
    pdfauthor={wuxch},
    unicode=false,          % xetex should set to false
    colorlinks,             % false: boxed links; true: colored links
    linkcolor=blue,         % color of internal links
    citecolor=green,        % color of links to bibliography
    filecolor=magenta,      % color of file links
    urlcolor=cyan           % color of external links
}
\makeindex

\usepackage[dvips,dvipsnames,svgnames]{xcolor}
\definecolor{light-gray}{gray}{0.95}

% 嵌入的代码显示
\usepackage{listings}
\lstset{language=C, breaklines, extendedchars=false}
\lstset{basicstyle=\ttfamily\small,
        tabsize=4,
        frame=single,
        framerule=0pt,
        backgroundcolor=\color{light-gray},
        % keywordstyle=\color{blue},
        keywordstyle=\bfseries,
        % commentstyle=\color{SeaGreen},
        stringstyle=\ttfamily,
        showstringspaces=false,
        escapebegin=\begin{color}{red}\begin{bfseries}\begin{large},
        escapeend=\end{large}\end{bfseries}\end{color}
        }

% \usepackage{graphicx}
% \usepackage{wrapfig}
% \usepackage{picinpar}

\renewcommand\contentsname{\center{目~~~~~录}}
\renewcommand\listfigurename{\center{插~~~~~图}}
\renewcommand\listtablename{表格}
\renewcommand\indexname{索引}
\renewcommand\figurename{图}
\renewcommand\tablename{表}


\setlength{\parskip}{0.8ex plus 0.2ex minus 0.2ex}
\setlength{\parindent}{0ex}

\usepackage[sf]{titlesec}
% \titleformat*{\section}{\large\sffamily}
% \renewcommand{\sectionname}{\CJKnumber{\thesection}}
\titleformat{\section}{\normalsize\sffamily}{\thesection}{1em}{}
\titleformat{\subsection}{\normalsize\sffamily}{\thesubsection}{1em}{}
\titlespacing*{\section}{0.5ex}{0.5ex}{0.5ex}
\titlespacing*{\subsection}{0.5ex}{0.2ex}{0.2ex}

\usepackage{fancybox}

\newcommand\blk[1]{\begin{minipage}{0.25\textwidth}#1 \rule[-1pt]{0.6\textwidth}{1pt}\end{minipage}}
\setlength\fboxrule{1.2pt}
\newcommand\score[1]{\hspace{2em}\color{blue}\framebox[3em][c]{\textbf\textit{{#1分}}}}
\newcommand\lookfordiff[1]{请找出下面代码中的隐患或者错误 (#1分)}

% \usepackage{fullpage}


% \newenvironment{mycode}%
% {\begin{footnotesize}%
% \begin{lstlisting}}%
% {\end{lstlisting}%
% \end{footnotesize}}

\usepackage{caption}
\renewcommand{\captionfont}{\scriptsize \sffamily}
\setlength{\abovecaptionskip}{0pt}
\setlength{\belowcaptionskip}{0pt}

\usepackage{fancyhdr}
\usepackage{lastpage}
\fancyhf{}


% \rhead{\scriptsize{\textsf{NGN分中心编程技能竞赛200906}}}
\cfoot{\scriptsize{\textsf{第 \thepage ~页，共 \pageref*{LastPage} 页}}}
\pagenumbering{arabic}
\pagestyle{fancy}

% ZTE标志，最后使用
\lhead{\includegraphics[height=2ex]{zte_logo.jpg}}
% 水印，最后使用
\usepackage{wallpaper}
\CenterWallPaper{1}{zte_confidential.png}

\usepackage{enumitem}
\setitemize{label=$\bullet$,leftmargin=3em,noitemsep,topsep=0pt,parsep=0pt}
\setenumerate{leftmargin=3em,noitemsep,topsep=0pt,parsep=0pt}
\renewcommand{\labelenumi}{(\arabic{enumi})}

% myblank
\newcommand{\myblank}[1]{\hspace{0.2em}\rule[-1pt]{#1}{0.3mm}\hspace{0.2em}}

% myoptions
\newcommand{\oneoption}[2]{%
\begin{minipage}[t]{.4\columnwidth}
\begin{minipage}[t]{2em}{%
#1)%
}%
\end{minipage}%
\begin{minipage}[t]{1\textwidth}{%
#2}%
\end{minipage}%
\end{minipage}%
}


\newcommand{\myoptions}[4]{%
\vspace{1ex}
\oneoption{A}{#1}\oneoption{B}{#2}

\vspace{1ex}
\oneoption{C}{#3}\oneoption{D}{#4}
\vspace{2ex}
}
